\chapter{Những câu hỏi về tương lai}
\markboth{Những câu hỏi về tương lai}{Những câu hỏi về tương lai}

\section{AI làm hết rồi, học làm gì?}
\begin{truyen}{AI làm hết rồi, học làm gì?}{Classroom Talk}
\chuhoa{M}{ột bạn giơ tay: ``Sau này AI làm hết rồi, học làm gì nữa thầy?''}

Thầy hỏi: ``Máy làm việc gì?''

``Dạ, làm việc nhanh.''

``Còn con người phải biết đặt câu hỏi, biết chọn điều đúng, biết chịu trách nhiệm. Máy không sống thay em được.''

Bạn ấy gật đầu, mặt bớt đắc thắng hơn lúc đầu.

\textit{(A student questions learning in the age of AI. The teacher distinguishes doing tasks from thinking, choosing, and taking responsibility.)}
\end{truyen}

\begin{baihoc}
Khi công cụ mạnh hơn, con người càng phải mạnh ở phần suy nghĩ và đạo đức.
\textit{(As tools grow stronger, people must grow stronger in thinking and ethics.)}
\end{baihoc}

\begin{ghinhoanh}
\textbf{Short English to remember:}

Tools can work fast.

Humans must think well.
\end{ghinhoanh}

\begin{tuhocnhanh}
1. Em đang muốn giao luôn phần suy nghĩ cho máy phải không? \textit{(Are you trying to hand thinking over to machines?)}

2. Phần nào chỉ con người mới chịu trách nhiệm được? \textit{(Which part can only humans take responsibility for?)}
\end{tuhocnhanh}
\ngancach

\section{Sau này em chưa biết làm nghề gì}
\begin{truyen}{Sau này em chưa biết làm nghề gì}{Classroom Talk}
\chuhoa{C}{ó bạn nói: ``Em chưa biết sau này làm nghề gì nên học không có động lực.''}

Thầy đáp: ``Chưa biết nghề thì càng phải học nền cho chắc.''

Bạn hỏi: ``Sao lại vậy?''

Thầy nói: ``Vì khi chưa biết đi đâu, điều tốt nhất là đừng tự làm yếu chân mình.''

Bạn ấy cười: ``Nghe hơi đau mà đúng.''

\textit{(A student lacks motivation because the future career is unclear. The teacher explains that uncertainty is exactly why a strong foundation matters.)}
\end{truyen}

\begin{baihoc}
Không biết tương lai không phải lý do để buông hiện tại.
\textit{(Not knowing the future is not a reason to drop the present.)}
\end{baihoc}

\begin{ghinhoanh}
\textbf{Short English to remember:}

Unclear future, stronger basics.

Do not weaken yourself early.
\end{ghinhoanh}

\begin{tuhocnhanh}
1. Em đang trì hoãn điều gì vì chưa rõ tương lai? \textit{(What are you delaying because the future is unclear?)}

2. Nền tảng nào em cần giữ chắc trước? \textit{(Which foundation do you need first?)}
\end{tuhocnhanh}
\ngancach

\section{Nhà em không khá thì mơ lớn được không?}
\begin{truyen}{Nhà em không khá thì mơ lớn được không?}{Classroom Talk}
\chuhoa{M}{ột bạn hỏi rất nhỏ: ``Nhà em không khá, em mơ lớn có vô lý không thầy?''}

Thầy đáp ngay: ``Không vô lý.''

Bạn nhìn lên.

Thầy nói tiếp: ``Điều kiện khó làm đường đi gập ghềnh hơn, chứ không cấm em có ước mơ. Nhưng mơ lớn thì phải đi thật bền.''

Bạn ấy ngồi rất yên, như sợ động là câu nói sẽ trôi mất.

\textit{(A student from a difficult family background wonders whether big dreams are unrealistic. The teacher answers with honesty: dreams remain possible, but the road requires greater endurance.)}
\end{truyen}

\begin{baihoc}
Hoàn cảnh có thể làm hành trình khó hơn, nhưng không có quyền cấm mình cố gắng.
\textit{(Circumstances may make the journey harder, but they do not own the right to forbid effort.)}
\end{baihoc}

\begin{ghinhoanh}
\textbf{Short English to remember:}

Hard beginnings do not cancel big dreams.

They demand stronger endurance.
\end{ghinhoanh}

\begin{tuhocnhanh}
1. Khó khăn nào đang làm em ngại mơ lớn? \textit{(What hardship makes you afraid to dream big?)}

2. Em cần sức bền ở điểm nào nhất? \textit{(Where do you need endurance the most?)}
\end{tuhocnhanh}
\ngancach

\section{Học mấy môn này ngoài đời để làm gì?}
\begin{truyen}{Học mấy môn này ngoài đời để làm gì?}{Classroom Talk}
\chuhoa{M}{ột bạn hỏi quen thuộc: ``Ngoài đời ai dùng hết mấy môn này đâu thầy?''}

Thầy đáp: ``Ngoài đời người ta dùng trí nhớ, cách suy luận, khả năng diễn đạt, sự kiên nhẫn và kỷ luật. Các môn học chỉ là nơi luyện những thứ đó.''

Bạn nói: ``Vậy môn học là cái cớ hả thầy?''

Thầy cười: ``Không phải cái cớ. Nó là phòng tập.''

Cả lớp bật cười vì ví dụ quá dễ hiểu.

\textit{(A student questions the practical use of subjects. The teacher explains that subjects are training grounds for broader abilities used in life.)}
\end{truyen}

\begin{baihoc}
Môn học không chỉ cho kiến thức. Nó còn rèn cách mình dùng đầu óc và tính cách.
\textit{(Subjects do not only give knowledge. They also train how we use mind and character.)}
\end{baihoc}

\begin{ghinhoanh}
\textbf{Short English to remember:}

Subjects are training rooms.

They build more than facts.
\end{ghinhoanh}

\begin{tuhocnhanh}
1. Một môn em ghét đang rèn cho em phẩm chất gì? \textit{(What quality is a disliked subject training in you?)}

2. Em có nhìn môn học rộng hơn nội dung bài chưa? \textit{(Do you see subjects beyond the lesson content?)}
\end{tuhocnhanh}
\ngancach

\section{Tương lai có chắc chắn không?}
\begin{truyen}{Tương lai có chắc chắn không?}{Classroom Talk}
\chuhoa{M}{ột bạn hỏi: ``Thầy ơi, cố gắng rồi tương lai có chắc chắn tốt hơn không?''}

Thầy lắc đầu: ``Không ai dám hứa chắc.''

Bạn hỏi tiếp: ``Vậy cố làm gì?''

Thầy đáp: ``Vì không cố thì gần như chắc chắn tệ hơn. Cố gắng không bảo đảm tất cả, nhưng nó nâng xác suất phần tốt về phía em.''

Bạn ấy cười: ``Nghe giống toán quá thầy.''

\textit{(A student wants certainty from effort. The teacher refuses false promises but explains that effort changes the odds in our favor.)}
\end{truyen}

\begin{baihoc}
Cố gắng không phải vé bảo đảm. Nó là cách làm cơ hội đứng gần mình hơn.
\textit{(Effort is not a guarantee ticket. It is a way to bring opportunity closer.)}
\end{baihoc}

\begin{ghinhoanh}
\textbf{Short English to remember:}

Effort gives no perfect guarantee.

It improves your chances.
\end{ghinhoanh}

\begin{tuhocnhanh}
1. Em có đang đòi sự chắc chắn rồi mới chịu cố không? \textit{(Are you demanding certainty before effort?)}

2. Điều gì đáng để em tăng xác suất thành công? \textit{(What is worth improving your chances for?)}
\end{tuhocnhanh}
\ngancach

\section{Nếu em chọn sai nghề thì sao?}
\begin{truyen}{Nếu em chọn sai nghề thì sao?}{Classroom Talk}
\chuhoa{M}{ột bạn hỏi lo lắng: ``Lỡ em chọn sai nghề thì sao thầy?''}

Thầy đáp: ``Ai cũng sợ điều đó.''

Bạn hỏi tiếp: ``Vậy làm sao đỡ sai?''

Thầy nói: ``Mình tìm hiểu kỹ hơn, trải nghiệm sớm hơn, và hiểu bản thân thật hơn. Còn nếu lỡ sai, quan trọng là dám điều chỉnh chứ không cố chấp giữ một hướng chỉ vì sợ mất mặt.''

Bạn ấy thở ra như bớt được một phần gánh nặng.

\textit{(A student fears choosing the wrong career. The teacher replaces panic with research, self-understanding, and the courage to adjust.)}
\end{truyen}

\begin{baihoc}
Chọn sai không đáng sợ bằng biết không hợp mà vẫn cố chấp đi tiếp chỉ vì sĩ diện.
\textit{(Choosing wrongly is less dangerous than stubbornly staying because of pride.)}
\end{baihoc}

\begin{ghinhoanh}
\textbf{Short English to remember:}

Wrong choices can be adjusted.

Pride can trap you longer.
\end{ghinhoanh}

\begin{tuhocnhanh}
1. Em đang sợ sai hay sợ bị nhìn là sai? \textit{(Are you afraid of being wrong, or of being seen as wrong?)}

2. Em cần tìm hiểu thêm điều gì về mình? \textit{(What do you need to learn more about yourself?)}
\end{tuhocnhanh}
\ngancach

\section{Muốn giàu nhanh có sai không?}
\begin{truyen}{Muốn giàu nhanh có sai không?}{Classroom Talk}
\chuhoa{C}{ó bạn hỏi thật: ``Muốn giàu nhanh có sai không thầy?''}

Thầy đáp: ``Muốn thì không sai.''

Bạn cười: ``Vậy tốt rồi.''

Thầy nói tiếp: ``Nhưng nếu chỉ muốn nhanh mà không muốn giỏi, không muốn bền, không muốn tử tế, thì cái nhanh đó thường rất đắt.''

Bạn ấy gật đầu, nụ cười nhỏ lại.

\textit{(A student asks whether wanting quick wealth is wrong. The teacher points not to the desire itself but to what is sacrificed for speed.)}
\end{truyen}

\begin{baihoc}
Nhanh không phải lúc nào cũng xấu. Nhưng nhanh mà rỗng thì rất dễ đổ.
\textit{(Fast is not always bad. But fast and hollow collapses easily.)}
\end{baihoc}

\begin{ghinhoanh}
\textbf{Short English to remember:}

Fast is tempting.

Strong and honest must come too.
\end{ghinhoanh}

\begin{tuhocnhanh}
1. Em đang thích cái nhanh ở điểm nào? \textit{(Where are you tempted by speed?)}

2. Em có đang bỏ qua phần nền tảng không? \textit{(Are you ignoring the foundation?)}
\end{tuhocnhanh}
\ngancach

\section{Em không có quan hệ thì sao?}
\begin{truyen}{Em không có quan hệ thì sao?}{Classroom Talk}
\chuhoa{M}{ột bạn than: ``Nhà em không quen biết ai, sau này em sao cạnh tranh lại người ta?''}

Thầy đáp: ``Quan hệ là một lợi thế, nhưng không phải tất cả.''

Bạn hỏi: ``Nhưng người có sẵn vẫn hơn mà thầy.''

Thầy gật đầu: ``Có, hơn ở một số cửa. Nhưng em vẫn còn những thứ phải tự xây: năng lực, uy tín, thái độ, cách làm việc. Có người được mở cửa sẵn mà không giữ nổi. Có người tự gõ cửa mà rồi người ta nhớ lâu.''

Bạn ấy nghe rất chăm chú.

\textit{(A student worries about lacking connections. The teacher acknowledges the unfairness but points back to what can still be built personally.)}
\end{truyen}

\begin{baihoc}
Mình không chọn được điểm xuất phát, nhưng vẫn chọn được cách mình tự làm mình đáng tin hơn.
\textit{(We do not choose our starting line, but we can still choose how trustworthy we become.)}
\end{baihoc}

\begin{ghinhoanh}
\textbf{Short English to remember:}

Connections help.

Character and ability still matter.
\end{ghinhoanh}

\begin{tuhocnhanh}
1. Em đang bi quan vì thiếu điều kiện nào? \textit{(What missing advantage is making you pessimistic?)}

2. Em còn chủ động xây được điều gì? \textit{(What can you still build on your own?)}
\end{tuhocnhanh}
\ngancach

\section{Có cần ước mơ lớn không?}
\begin{truyen}{Có cần ước mơ lớn không?}{Classroom Talk}
\chuhoa{M}{ột bạn hỏi: ``Ai cũng nói phải có ước mơ lớn. Nếu em chỉ muốn sống bình an thôi thì có kém quá không?''}

Thầy mỉm cười: ``Không kém.''

Bạn ngước lên.

Thầy nói tiếp: ``Ước mơ không cần phải ồn ào mới có giá trị. Chỉ cần nó đủ tử tế và đủ thật với em. Điều quan trọng không phải mơ lớn cỡ nào, mà là em sống với điều mình chọn có nghiêm túc không.''

Bạn ấy cười nhẹ, có vẻ được giải tỏa.

\textit{(A student wonders if a simple dream is inferior. The teacher gives dignity to quiet aspirations lived sincerely.)}
\end{truyen}

\begin{baihoc}
Không phải ai cũng phải sống như tiêu đề lớn. Nhiều cuộc đời tử tế được xây rất lặng lẽ.
\textit{(Not everyone must live like a grand headline. Many good lives are built quietly.)}
\end{baihoc}

\begin{ghinhoanh}
\textbf{Short English to remember:}

Not every dream must be loud.

Quiet lives can be meaningful.
\end{ghinhoanh}

\begin{tuhocnhanh}
1. Em có đang xấu hổ vì ước mơ của mình quá giản dị không? \textit{(Are you ashamed that your dream feels too simple?)}

2. Em muốn sống nghiêm túc với điều gì? \textit{(What do you want to live seriously for?)}
\end{tuhocnhanh}
\ngancach

\section{Nếu em chưa có ước mơ thì sao?}
\begin{truyen}{Nếu em chưa có ước mơ thì sao?}{Classroom Talk}
\chuhoa{C}{ó bạn thở dài: ``Người ta hỏi ước mơ hoài mà em chưa có. Vậy em có bị chậm không thầy?''}

Thầy đáp: ``Chưa có chưa chắc là chậm.''

Bạn hỏi: ``Vậy là gì?''

Thầy nói: ``Có thể em đang trong giai đoạn đi tìm. Nhưng trong lúc tìm, em vẫn phải sống cho đàng hoàng và học cho chắc. Không có ước mơ rõ ràng không đồng nghĩa với việc được phép buông thả.''

Bạn ấy cười ngượng nhưng gật đầu rất thật.

\textit{(A student has no clear dream yet. The teacher normalizes that uncertainty while refusing to let it become an excuse for drifting.)}
\end{truyen}

\begin{baihoc}
Chưa có câu trả lời lớn không sao. Miễn là mình vẫn giữ những bước nhỏ cho đàng hoàng.
\textit{(It is okay not to have the big answer yet, as long as you still keep your small steps in order.)}
\end{baihoc}

\begin{ghinhoanh}
\textbf{Short English to remember:}

No clear dream is not failure.

Keep your small steps steady.
\end{ghinhoanh}

\begin{tuhocnhanh}
1. Em đang lo vì chưa biết đường hay vì sợ bị hỏi? \textit{(Are you worried because you do not know the path, or because people keep asking?)}

2. Bước nhỏ nào em vẫn phải giữ? \textit{(What small step must you still keep?)}
\end{tuhocnhanh}
\ngancach